% UNBank design document. Compiles with pdflatex on Overleaf or with latexmk -pdf main.tex.
% Mermaid diagrams go in figures/ as PDF (mmdc -i x.mmd -o x.pdf) and replace the
% \mermaidplaceholder boxes. Remove every \plan{...} box before submission.
\documentclass[11pt,a4paper]{article}

\usepackage[T1]{fontenc}
\usepackage[utf8]{inputenc}
\usepackage{lmodern}
\usepackage{microtype}
\usepackage[a4paper,margin=2.5cm]{geometry}
\usepackage{graphicx}
\usepackage{booktabs}
\usepackage{tabularx}
\usepackage{longtable}
\usepackage{array}
\usepackage{xcolor}
\usepackage{float}
\usepackage{caption}
\usepackage{enumitem}
\usepackage{titlesec}
\usepackage{fancyhdr}
\usepackage[hidelinks]{hyperref}
\usepackage[nameinlink,capitalise]{cleveref}

% Ragged-right table columns, so narrow cells don't stretch their word spacing.
\newcolumntype{L}[1]{>{\raggedright\arraybackslash}p{#1}}
\newcolumntype{Y}{>{\raggedright\arraybackslash}X}

\definecolor{planbg}{HTML}{EEF2FB}
\definecolor{planrule}{HTML}{2D2AE6}
\definecolor{figbg}{HTML}{F5F5F7}

\newcommand{\project}{Unbeatable Bank}
\newcommand{\repourl}{\url{https://github.com/MSRishab007/unbeatable-bank}}
\newcommand{\liveurl}{\url{https://unbeatable-bank.duckdns.org}}
\newcommand{\code}[1]{\texttt{#1}}

% Drafting aid: what a section will contain. Delete before submission.
\newcommand{\plan}[1]{%
  \par\noindent\fcolorbox{planrule}{planbg}{\parbox{\dimexpr\linewidth-2\fboxsep-2\fboxrule}{\small\textbf{Section plan.} #1}}\par\medskip}

% Placeholder for a Mermaid diagram exported to PDF later.
% #1 label, #2 caption, #3 what the diagram shows, #4 height
\newcommand{\mermaidplaceholder}[4]{%
  \begin{figure}[H]\centering
  \fcolorbox{gray}{figbg}{\parbox[c][#4][c]{0.9\linewidth}{\centering\small\textit{Mermaid diagram to be inserted.}\\[4pt]#3\\[4pt]\texttt{figures/#1.pdf}}}
  \caption{#2}\label{fig:#1}
  \end{figure}}

% Placeholder for a screenshot.
\newcommand{\screenshotplaceholder}[4]{%
  \begin{figure}[H]\centering
  \fcolorbox{gray}{white}{\parbox[c][#4][c]{0.9\linewidth}{\centering\small\textit{Screenshot to be inserted.}\\[4pt]#3\\[4pt]\texttt{figures/#1.png}}}
  \caption{#2}\label{fig:#1}
  \end{figure}}

\pagestyle{fancy}
\fancyhf{}
\fancyhead[L]{\small \project{}}
\fancyhead[R]{\small Design Document}
\fancyfoot[C]{\small \thepage}

\setlist[itemize]{leftmargin=1.5em,itemsep=2pt}
\captionsetup{font=small,labelfont=bf}

\begin{document}

% =====================================================================  title page
\begin{titlepage}
\centering
{\large Birla Institute of Technology and Science, Pilani\par}
\vspace{2.5cm}
{\Huge\bfseries \project{}\par}
\vspace{0.4cm}
{\Large Design of a Cloud-Native Retail Banking Platform\\for Transaction Monitoring\par}
\vspace{0.6cm}
{\large Design Document\par}
\vspace{1.5cm}
{\normalsize Scenario: Banking Transaction Monitoring\\Concept: Leaderless Replication and Dynamo-Style Storage\par}
\vspace{2cm}
\begin{tabular}{ll}
\toprule
\textbf{ID number} & \textbf{Name} \\
\midrule
20XXXXXXXXXX & Member 1 \\
20XXXXXXXXXX & Member 2 \\
20XXXXXXXXXX & Member 3 \\
20XXXXXXXXXX & Member 4 \\
\bottomrule
\end{tabular}
\vfill
{\small Source repository: \repourl\\[2pt]
Live deployment (available during the demonstration): \liveurl\par}
\vspace{0.5cm}
{\small November 2026\par}
\end{titlepage}

% =====================================================================  abstract
\begin{abstract}
\noindent Payment systems must admit legitimate transactions within milliseconds while they detect fraud that only emerges across many transactions. This document presents the design of \project{}, a cloud-native retail banking platform for transaction admission and fraud monitoring. The platform consists of twelve microservices on Google Cloud. A managed Kubernetes cluster hosts the stateless services, a serverless function delivers notifications, and a Flink job on Dataproc analyses the transaction stream. The services communicate through REST, gRPC, Kafka and Pub/Sub. The design separates two kinds of data. Balances and journal entries need strong consistency, so they live in PostgreSQL behind a single Ledger service. Rate-limit counters and risk features need high availability, so they live in DynaStore, a leaderless key-value store that implements the techniques of Amazon's Dynamo. Each payment passes five layers of defence. These are admission control, authorisation and policy, real-time risk scoring with 21 weighted rules, stream pattern detection, and case management. Terraform provisions all infrastructure, and ArgoCD deploys every release from a Git repository. In a sustained load test, the platform scaled out automatically, failed 0.01\% of requests and answered 95\% of requests within 1.4 seconds. A fraud simulator showed that the layers stopped all seven attack scenarios without disturbing an ordinary customer.
\end{abstract}

\tableofcontents
\newpage

% =====================================================================  1
\section{Introduction}
\label{sec:introduction}

\subsection{Scenario and Objectives}
\label{sec:scenario}

A retail bank processes a continuous stream of payments, and a small share of them is fraudulent or abusive. The bank must decide within a fraction of a second whether to admit each payment. It must also watch the stream over minutes and hours for patterns that no single payment reveals. A collection account that receives money from many unrelated customers is one such pattern. A series of payments kept just below a reporting threshold is another.

This document presents \project{}, a cloud-native retail banking platform built for scenario S17, Banking Transaction Monitoring. The scenario covers transaction admission and fraud monitoring. It pairs this domain with the concept of leaderless replication in a storage system modelled on Amazon's Dynamo~\cite{dynamo}.

The scenario suggests eight microservices. These are the Account, Transaction Intake, Rate Limit Coordinator, Ledger, Risk Detection, Fraud Alert, Notification and Audit Analytics services. The platform implements all eight of them. It adds four services that a working bank needs. They are an API Gateway, an Identity and Authentication service, a Deposits service, and DynaStore, a purpose-built Dynamo-style key-value store. The platform runs on Google Cloud. Customers use it through a public web application, and bank staff use a console within the same application.

The platform follows current practice for cloud-native systems. Code provisions all of its infrastructure. A serverless function performs asynchronous work, and the services communicate through REST, gRPC and message queues. A managed Kubernetes service hosts the stateless services and scales them with horizontal pod autoscalers. A GitOps controller performs every deployment. A stream processing job on a managed cluster analyses the transaction stream. The data sits in several cloud storage products, and a monitoring and logging stack observes the whole system. \Cref{sec:traceability} maps each of these properties to the part of the system that provides it.

The project pursues three further objectives. First, every payment passes several independent checks before any money moves. Second, the system continues to answer when individual components fail. Third, every change to the infrastructure or the software reaches the cloud through version control, and no change is made by hand.

\subsection{Guiding Design Principles}
\label{sec:principles}

Four principles guide the design. The first two concern data, and they lead to the most important decision in the system.

\paragraph{Consistency where money moves.} Money must never be created, lost or moved twice. This property requires strong consistency. The system therefore keeps balances and journal entries in a relational database, and only the Ledger service may change them. Each debit runs in a single database transaction, and database constraints reject any journal whose debits and credits do not balance.

\paragraph{Availability where admission is decided.} Admission control must answer quickly on every request, including when parts of the system fail. The rate limiter and the risk engine update counters and features on every request. A value that is a few milliseconds old does no harm to these decisions. Such data suits leaderless replication, in which every replica accepts reads and writes and a request succeeds once a quorum of replicas has answered. The system keeps this data in DynaStore. The separation of the two kinds of data, each in a store with the guarantees it needs, is the central design decision of \project{}.

\paragraph{Defence in depth.} No single check can stop every attack. Each payment therefore passes five layers, which are admission control, authorisation and policy, real-time risk scoring, stream analytics and case management. A payment that one layer misses can still be caught by another, and the slower layers feed what they learn back into the faster ones.

\paragraph{Everything declared in Git.} Terraform describes all cloud resources, and Kubernetes manifests describe all workloads. A release happens when a pull request that changes the image versions is merged. ArgoCD then brings the cluster in line with the repository. The repository therefore records the complete state of the system at every point in its history.

\subsection{Contributions}
\label{sec:contributions}

This project makes the following contributions.

\begin{itemize}
\item A layered model for payment admission. It combines distributed rate limiting, identity-based controls, a real-time risk engine with 21 weighted rules, streaming pattern detection, and case management with a maker-checker rule.
\item DynaStore, a leaderless key-value store that reproduces the main techniques of Dynamo. These include consistent hashing, tunable quorums, vector clocks, conflict-free replicated counters, hinted handoff, read repair and anti-entropy.
\item A Flink SQL job on Dataproc that computes per-minute transaction statistics, including the number of unique customers. The job also detects four fraud patterns and feeds them back into the real-time risk engine.
\item A fully automated delivery process. Terraform provisions all infrastructure, GitHub Actions tests the code and builds the images, and ArgoCD deploys every release from the repository.
\item An evaluation with three load test runs and a fraud simulator of eight scenarios, seven of which are attacks. In the final load test, 0.01\% of requests failed and 95\% of requests completed within 1.4 seconds. In the simulator, at least one layer stopped every attack scenario, while an ordinary customer passed without friction.
\end{itemize}

\subsection{Artefact Availability}
\label{sec:artefacts}

The source repository is available at \repourl{}. It contains the code of all twelve services, the Terraform configuration, the Kubernetes manifests, the ArgoCD configuration and the automated tests. The file \code{task\_division.txt} in the repository lists the team members and the tasks assigned to each of them.

The deployed system is reachable at \liveurl{} during demonstrations and test runs. Outside these periods, the cluster scales to zero nodes and the database stops, which keeps the cost within the available cloud credit. The address may therefore not respond at other times. The deployment guide in the repository describes how to recreate the full environment in a new Google Cloud project.

\subsection{Document Organisation}
\label{sec:organisation}

The rest of this document is organised as follows. \Cref{sec:overview} gives an overview of the system and maps its requirements to the implementation. \Cref{sec:deployment} describes the cloud deployment architecture, including the infrastructure code and the delivery pipeline. \Cref{sec:microservices} describes the microservices and their responsibilities, and \cref{sec:interconnection} describes how they communicate. \Cref{sec:fraud} presents the techniques for transaction admission and fraud detection in depth. \Cref{sec:data} describes the data architecture, including DynaStore and the stream processing job. \Cref{sec:quality} evaluates scalability, resilience, observability and security. \Cref{sec:rationale} collects the rationale for the main design choices, and \cref{sec:conclusion} discusses limitations and concludes the document.

% =====================================================================  2
\section{System Overview}
\label{sec:overview}

This section describes what \project{} does for its users and how it defends each payment. It then relates the services to those of scenario S17 and shows where each platform requirement is met.

\subsection{Functional Scope}
\label{sec:scope}

\project{} serves two groups of users, customers and bank staff. Both groups use the same web application, which the API Gateway serves from a single public address.

A customer registers with a name, date of birth, mobile number, email address and password. The customer then uploads identity documents for the know-your-customer (KYC) check. In this check, the bank verifies the identity of a new customer before money can move. Once the bank accepts these documents, the customer can hold any number of savings and current accounts, alone or jointly with another customer. The customer can also book fixed deposits and recurring deposits. A fixed deposit locks a sum for a chosen term at a fixed rate. A recurring deposit collects a fixed instalment every month. The application shows a quote with the maturity amount before the customer commits, and it prices any early closure in advance.

Customers pay one another by account number, by a chosen pay identifier such as \code{riya@unbeatable}, or by scanning a QR code. Every payment shows the name of the payee before any money moves. A new payee can receive only a limited amount during the first 24 hours. This delay gives the customer time to notice a payee that someone else added. After each payment, the application shows the checks that the payment passed and the check that stopped it, if any. Customers receive receipts, one-time passwords and security alerts in an in-app inbox, by email, and through a simulated RCS channel. They can download monthly statements, and anyone can verify that a statement file was issued by the bank.

Customers also control the security of their own accounts. Each customer sets a six-digit payment PIN for money transfers and a four-digit app PIN that unlocks the application. The application locks itself after five minutes without use. A customer can review the devices and sessions that hold access and sign any of them out. A customer can also cancel a payment that the bank has held for review and report a payment as fraudulent. In an emergency, a single action blocks all access to every account.

Bank staff work in a console that shows only the functions of their role. KYC officers review identity documents and accept or reject them. Fraud analysts work through alerts and cases, release or cancel held payments, and freeze accounts. Tellers record cash deposits at the branch desk. Administrators configure the notification channels and monitor the stream analytics and DynaStore. Sensitive actions follow a maker-checker rule, under which one staff member proposes the action and a second staff member approves it. Unfreezing an account is one such action. \Cref{fig:app-overview} shows the customer application.

\screenshotplaceholder{app-overview}{Customer application showing the home screen and the payment flow.}{Phone screenshots of Home, Pay with the swipe control, and the payment result with its six checks.}{5cm}

\subsection{Layered Payment Defence}
\label{sec:layers}

Every payment passes five layers of defence before and after money moves. \Cref{tab:layers} lists the layers in the order in which they act. The first three layers act synchronously, while the customer waits for the answer. The last two layers act on the stream of events that completed payments produce.

\begin{table}[H]\centering\small
\caption{Layers of defence applied to each payment.}\label{tab:layers}
\begin{tabularx}{\linewidth}{@{}L{2.4cm}L{2.3cm}L{2.9cm}Y@{}}
\toprule
Layer & Time scale & Service & Technique \\
\midrule
Admission & Before processing & Rate Limit Coordinator & Sliding-window limits per customer, account, device and IP address, kept as conflict-free counters in DynaStore \\
Authorisation and policy & Synchronous & API Gateway, Account, Transaction Intake & PIN check, holder and account status checks, and limits per payment, per day and per new payee \\
Real-time risk & Tens of milliseconds & Risk Detection & Twenty-one weighted rules over features in DynaStore, with a score that allows, challenges, holds or blocks the payment \\
Stream analytics & Seconds to minutes & Audit Analytics (Flink) & Windowed aggregations that detect bursts, probing, structuring and mule fan-in \\
Response & Minutes & Fraud Alert & Alerts, automatic freezes, cases with maker-checker review, and signals fed back to the risk engine \\
\bottomrule
\end{tabularx}
\end{table}

The admission layer stops abusive volumes cheaply, before any other service does work. The authorisation layer confirms that the customer may use the account and that the payment respects the account limits. The risk layer scores the payment against the recent behaviour of the customer, the device and the payee. A payment with a moderate score needs a one-time password, and a payment with a high score waits for an analyst. The stream analytics layer looks at many payments together and finds patterns that no single payment reveals. The response layer turns these findings into alerts and cases. It also writes short-lived signals into DynaStore, so that the risk engine scores the next payment from a flagged account higher. The slow layers therefore improve the decisions of the fast layers over time. \Cref{sec:fraud} describes each layer in depth.

\subsection{Mapping to the Scenario Services}
\label{sec:s17map}

Scenario S17 suggests eight services. \Cref{tab:s17} lists each of them with its role in the scenario and the \project{} service that implements it. The platform keeps the names and the responsibilities of all eight services. The Notification Service runs as a serverless function and the Audit Analytics Service runs as a stream processing job. The other six services run on Kubernetes.

\begin{table}[H]\centering\small
\caption{Services suggested by scenario S17 and their implementation in \project{}.}\label{tab:s17}
\begin{tabularx}{\linewidth}{@{}L{3.6cm}YL{3.6cm}@{}}
\toprule
Scenario service & Role in the scenario & \project{} service \\
\midrule
Account Service & Maintains account metadata & Account \\
Transaction Intake Service & Accepts transaction requests & Transaction Intake \\
Rate Limit Coordinator & Keeps shared rate-limit state and allocates permits & Rate Limit Coordinator \\
Ledger Service & Records accepted transactions & Ledger \\
Risk Detection Service & Evaluates transactions against risk rules & Risk Detection \\
Fraud Alert Service & Reports suspicious transaction patterns & Fraud Alert \\
Notification Service & Sends transaction outcomes & Notification (Cloud Function) \\
Audit Analytics Service & Measures rate-limit decisions, throughput and rejected requests & Audit Analytics (Flink on Dataproc) \\
\bottomrule
\end{tabularx}
\end{table}

The platform adds four services to the eight. The API Gateway is the only public entry point. It verifies every request, checks the role of the caller and the ownership of the account, and serves the web application. The Identity and Authentication service manages registration, sign-in, sessions, devices, PINs, one-time passwords and KYC documents. Without it, the system could not tell a customer from an attacker who holds a stolen password. The Deposits service manages fixed and recurring deposits and runs the end-of-day batch, which accrues interest, collects instalments and pays out matured deposits. DynaStore implements the leaderless replication concept of the scenario. It holds the rate-limit counters and risk features that the Rate Limit Coordinator and the Risk Detection service read and write on every payment.

\subsection{Requirement Traceability}
\label{sec:traceability}

\Cref{tab:traceability} maps each platform requirement to the part of the system that meets it and to the section that discusses it in depth.

\begin{table}[htbp]\centering\small
\caption{Mapping of platform requirements to the implementation.}\label{tab:traceability}
\begin{tabularx}{\linewidth}{@{}L{3.6cm}YL{1.7cm}@{}}
\toprule
Requirement & Implementation & Section \\
\midrule
Public web service and at least eight microservices & Twelve services, with the API Gateway reachable over HTTPS through a Google load balancer and Traefik & \cref{sec:microservices} \\
Infrastructure as Code & Terraform in three layers provisions the network, cluster, databases, storage, IAM policies, secrets, function and Dataproc cluster & \cref{sec:iac} \\
Serverless function & The Notification Service runs as a Cloud Function that Pub/Sub triggers asynchronously & \cref{sec:interconnection} \\
REST, gRPC and message queues & REST between the client and the services, gRPC on the payment path and inside DynaStore, Kafka and Pub/Sub for events & \cref{sec:interconnection} \\
Managed Kubernetes, autoscaling and load test & GKE hosts the stateless services, six horizontal pod autoscalers scale on CPU, and k6 generates sustained load & \cref{sec:scalability} \\
GitOps controller & ArgoCD tracks the repository and applies every release, and no deployment uses kubectl apply & \cref{sec:gitops} \\
Stream processing on a managed cluster & A Flink SQL job on Dataproc reads a Kafka topic, aggregates over time windows and writes results to separate Kafka topics & \cref{sec:streaming} \\
Leaderless replication concept & DynaStore implements consistent hashing, quorums, vector clocks, conflict-free counters, hinted handoff, read repair and anti-entropy & \cref{sec:concept} \\
Multiple storage products & Cloud SQL for PostgreSQL, Firestore and Cloud Storage, together with DynaStore and Kafka & \cref{sec:storage} \\
Observability & Prometheus and Grafana for metrics, and Cloud Logging for the logs of every service and of the Flink job & \cref{sec:observability} \\
\bottomrule
\end{tabularx}
\end{table}

% =====================================================================  3
\section{Cloud Deployment Architecture}
\label{sec:deployment}

This section describes how \project{} is deployed on Google Cloud. It covers the compute platform, the network and the public edge, the managed services, the infrastructure code, the delivery pipeline, access control and cost management.

\subsection{Architecture Overview}
\label{sec:arch-overview}

\Cref{fig:deployment} shows the deployment. All resources belong to one Google Cloud project in the \code{us-central1} region and share one Virtual Private Cloud (VPC) network. A Google Kubernetes Engine (GKE) cluster runs the stateless services, DynaStore, Kafka and the platform components. These platform components are the edge proxy, the GitOps controller and the monitoring stack. Managed services outside the cluster hold the relational, document and object data. They also run the notification function and the stream processing job.

A request from a customer travels as follows. The browser resolves the public host name to the address of a Google load balancer. The load balancer forwards the connection to Traefik, an edge proxy inside the cluster that terminates HTTPS. Traefik passes the request to the API Gateway, which authenticates it and calls the internal services. These services read and write Cloud SQL, Firestore, Cloud Storage and DynaStore. They publish events to Kafka and notification requests to Pub/Sub. The Flink job on Dataproc consumes the Kafka events, and the Cloud Function consumes the Pub/Sub messages. Cloud Logging collects the logs of every component, and Prometheus collects the metrics of every service.

\mermaidplaceholder{deployment}{Cloud deployment architecture of \project{} on Google Cloud.}{Users and the load test reach a Google load balancer, then Traefik and the API Gateway in GKE. The bank namespace holds the services, DynaStore and Kafka. Cloud SQL, Firestore, Cloud Storage, Pub/Sub, the Cloud Function, Dataproc, Secret Manager, Artifact Registry and Cloud Logging sit around the cluster inside one VPC.}{7cm}

\subsection{Compute Platform}
\label{sec:compute}

GKE hosts the containerised workloads. The cluster is zonal and runs in \code{us-central1-a}. It follows the regular release channel, so Google applies tested Kubernetes upgrades automatically. The cluster uses GKE Dataplane V2, which enforces Kubernetes network policies in the data path.

One node pool provides the compute capacity. Its nodes are \code{e2-standard-4} machines with four virtual CPUs and 16\,GB of memory each. The pool normally runs two nodes, and an operator can raise the count through Terraform when a load test needs more capacity. The nodes are Spot virtual machines. Google offers Spot capacity at a large discount, but it may reclaim a Spot machine at short notice. The workloads tolerate such a loss. The services keep no state in their pods, so Kubernetes simply restarts them on the remaining node. DynaStore and Kafka keep their data on persistent disks, which survive the loss of a node.

Workloads live in separate namespaces. The \code{bank} namespace holds the twelve services, DynaStore and the scheduled jobs. Other namespaces hold Kafka, Traefik, cert-manager, ArgoCD and the monitoring stack. This separation lets each component carry its own access rules and network policies.

\subsection{Networking and Edge}
\label{sec:network}

The VPC has one subnet for the nodes and two secondary ranges for pods and Kubernetes services. A firewall rule admits traffic between these internal ranges. A second rule admits SSH only from the address range of Google's Identity-Aware Proxy. Operators can therefore reach the Dataproc machine without exposing it to the internet.

Cloud SQL has only a private IP address. The VPC reaches it through Private Service Access, a peering between the VPC and the network that Google manages for the database. No address outside the VPC can connect to the database.

The public edge consists of a Google load balancer and Traefik. Kubernetes creates the load balancer for the Traefik service, which keeps the client address intact on its way to the pods. Traefik runs as two to four replicas under an autoscaler and a disruption budget. It terminates HTTPS with a certificate that cert-manager obtains from Let's Encrypt. Traefik redirects every plain HTTP request to HTTPS and sends a Strict-Transport-Security header, which tells browsers to use HTTPS for all later visits. Two DuckDNS host names point at the load balancer. The public address of the load balancer changes when the environment is recreated. A scheduled job therefore checks the address every two minutes and updates both names.

Inside the cluster, a network policy in the \code{bank} namespace admits traffic only from the same namespace and from the edge proxy. The API Gateway is therefore the only service that external traffic can reach. The internal services accept requests from the gateway and from one another.

\subsection{Managed Cloud Services}
\label{sec:managed}

\Cref{tab:managed} lists the managed services and the purpose of each one. \Cref{sec:data} explains why each kind of data lives in its particular store.

\begin{table}[htbp]\centering\small
\caption{Managed Google Cloud services used by \project{}.}\label{tab:managed}
\begin{tabularx}{\linewidth}{@{}L{3.4cm}Y@{}}
\toprule
Service & Purpose in the system \\
\midrule
GKE & Hosts the stateless services, DynaStore, Kafka and the platform components \\
Cloud SQL for PostgreSQL & Holds customers, accounts, the ledger, payments, deposits, credentials and cases, with one database user and schema for each service \\
Firestore & Holds sessions, login events, alerts, the notification inbox and analytics snapshots \\
Cloud Storage & Holds KYC documents, statements, end-of-day reports, Flink files and checkpoints, function source and Terraform state \\
Pub/Sub & Carries notification requests to the serverless function \\
Cloud Functions & Runs the Notification Service, which scales from zero to three instances \\
Dataproc & Runs the Flink stream processing job on a single-node cluster \\
Secret Manager & Stores database passwords, signing keys and provider credentials \\
Artifact Registry & Stores the container images that the delivery pipeline builds \\
Cloud Logging & Collects the logs of every container and of the Dataproc job \\
\bottomrule
\end{tabularx}
\end{table}

\subsection{Infrastructure as Code}
\label{sec:iac}

Terraform provisions every cloud resource of the system, and no resource is created by hand in the console. Terraform describes the desired state of the resources in configuration files and changes the cloud to match. The configuration is split into three layers, which \cref{fig:terraform-layers} shows. Each layer has its own state and depends on the outputs of the layer below it.

The bootstrap layer creates the foundation. It creates the project, enables the required APIs and grants the team members their roles. It also creates the bucket that stores the Terraform state of the other layers. It sets up Workload Identity Federation for GitHub Actions and creates budget alerts for each billing account. The main layer creates the workload infrastructure. It contains the network, the GKE cluster, Cloud SQL, Firestore and the storage buckets. It also contains the Pub/Sub topic, the Cloud Function, the Dataproc cluster, the service accounts and the secrets. The platform layer works inside the cluster. It creates the \code{bank} namespace, copies the secrets into Kubernetes, and installs ArgoCD together with its root application.

Terraform also generates the secrets. It creates a random password for each database user and a random key for signing tokens, and it stores them in Secret Manager. The configuration therefore never contains a secret, and the repository holds only example variable files. The Terraform state lives in a versioned Cloud Storage bucket, so every operator works from the same state.

\mermaidplaceholder{terraform-layers}{Layers of the Terraform configuration and the resources each layer manages.}{Three stacked layers, bootstrap, main and platform, with the resources each one creates and the dependencies between them.}{5cm}

\subsection{Continuous Delivery and GitOps}
\label{sec:gitops}

GitOps is a deployment practice in which a Git repository holds the desired state of the system. A controller then keeps the running system equal to that state. \project{} follows this practice with ArgoCD, and nothing deploys a service with \code{kubectl apply}. \Cref{fig:gitops} shows the pipeline from a commit to a running pod.

Every push to the main branch starts the continuous integration workflow in GitHub Actions. The workflow runs the unit tests of every service, the tests of the serverless function and the Flink SQL tests. It also validates the Kubernetes manifests, checks the Terraform configuration and builds the web application. When a push changes service code, a release workflow builds a container image for every service. It tags each image with the commit identifier and pushes it to Artifact Registry. The workflow authenticates to Google Cloud through Workload Identity Federation, so it holds no long-lived key. The workflow then opens a deployment pull request that changes only the image tags in the Kustomize configuration of the \code{bank} namespace.

Merging the deployment pull request is the release. ArgoCD tracks the main branch and detects the change within three minutes. It first runs the database migrations of each service as a pre-sync job. It then rolls out the new pods one at a time, so the service stays available during the update. ArgoCD runs with automatic synchronisation, self-healing and pruning. Self-healing reverts any manual change to the cluster, and pruning removes any resource that the repository no longer declares.

ArgoCD manages nine applications through the app-of-apps pattern. In this pattern, a root application declares the other applications as its resources. The root application declares the bank services, the Kafka operator and cluster, the monitoring stack and its dashboards, cert-manager, Traefik and the edge configuration. Third-party components such as Traefik and the monitoring stack come from Helm charts with pinned versions. The bank services come from plain manifests that a script generates from one description of all services.

\mermaidplaceholder{gitops}{Delivery pipeline from a commit to the cluster through GitHub Actions and ArgoCD.}{A commit to main triggers CI and the image build, which push to Artifact Registry and open a deploy pull request. Merging it updates the manifests on main, and ArgoCD syncs the cluster after running migrations.}{5cm}

\subsection{Identity and Access Management}
\label{sec:iam}

Each component receives only the permissions it needs. The system defines separate Google service accounts for the cluster nodes, the application services and the Identity service. Further accounts serve the notification function, its build, the Dataproc cluster and the delivery pipeline.

Pods authenticate to Google APIs through Workload Identity. This GKE feature maps a Kubernetes service account to a Google service account, so pods obtain short-lived credentials without any key file. Only three services call Google APIs directly, and only these three receive such a mapping. The Identity service has its own Google service account, and only this account may read the bucket that holds KYC documents.

Access to the database follows the same principle. Each service connects as its own database user and owns its own schema. A fault or a compromise in one service therefore cannot read or change the tables of another service. The containers run as a non-root user.

\subsection{Cost Management}
\label{sec:cost}

The design keeps the running cost within the available cloud credit. GKE charges no management fee for one zonal cluster, and Spot nodes cost a fraction of regular nodes. The database runs on a small shared-core tier. The notification function scales to zero instances when no messages arrive. Dataproc runs only during stream processing tests and demonstrations.

A single command stops the environment between working sessions. It scales the node pool to zero, stops the Cloud SQL instance and deletes the Dataproc cluster, all through Terraform. A second command restores the environment in a few minutes, and ArgoCD then brings every workload back from the repository. Budget alerts notify the team by email when spending reaches 30, 60 and 80 percent of the budget of each billing account.

% =====================================================================  4
\section{Microservices Architecture}
\label{sec:microservices}

This section describes how \project{} is divided into services, what each service is responsible for, and how the services cooperate to process a payment. It closes with the client applications that customers and staff use.

\subsection{Service Decomposition}
\label{sec:decomposition}

\project{} consists of twelve services. \Cref{fig:services} shows them with the calls between them. The decomposition follows business capabilities. Each service owns one capability of the bank, such as accounts, money movement or risk scoring. Each service also owns its data, and no other service reads or writes that data directly. A service that needs information from another service asks for it through the API of that service or learns it from an event.

This ownership has two benefits. First, a team can change the internal design of a service without breaking the others, as long as its API stays compatible. Second, each service can scale and fail independently. A slow risk engine, for example, cannot corrupt the ledger, because the ledger accepts changes only through its own interface.

Ten of the services run as containers on GKE. They are written in Python 3.12. They offer REST interfaces built with FastAPI, and four of them also offer gRPC interfaces. Every service exposes three standard endpoints. The \code{/healthz} endpoint reports that the process is alive. The \code{/readyz} endpoint reports that the service can accept traffic. The \code{/metrics} endpoint publishes Prometheus metrics. The Notification Service runs as a Cloud Function, and the Audit Analytics Service runs as a Flink job on Dataproc.

\mermaidplaceholder{services}{Microservices of \project{} and the interactions between them.}{The gateway calls the services over REST. Intake calls Rate Limit, Risk and Ledger over gRPC. Rate Limit and Risk use DynaStore. Kafka carries events to Flink and Fraud Alert, and Pub/Sub carries notifications to the Cloud Function.}{8cm}

\subsection{Service Responsibilities}
\label{sec:responsibilities}

This subsection describes each service in turn, grouped by area. \Cref{tab:services} at the end of the subsection summarises the responsibilities, the runtime and the data of every service.

\subsubsection{Edge and Identity}

\paragraph{API Gateway.} The API Gateway is the only service that the internet can reach. It verifies the token of every request and asks the Identity service whether the session behind the token is still active. It then checks the role of the caller and, for account operations, whether the caller holds the account. For sensitive operations, it verifies the payment PIN or the app PIN with the Identity service before it forwards the request. The gateway removes any identity header that a client sends and sets the header itself, so internal services can trust it. It also serves both web applications and limits the size of request bodies.

\paragraph{Identity and Authentication.} The Identity service manages everything that proves who a person is. It handles registration, sign-in, password changes and password resets. It issues short-lived access tokens and refresh tokens, and it rotates a refresh token every time the client uses it. It tracks the devices and sessions of each customer and locks a session after five minutes without activity. It stores the two PINs of each customer and locks a PIN after three wrong attempts. It issues and verifies the one-time passwords that confirm risky payments and new devices. It detects sign-ins from new devices and from implausibly distant places, and it reports both as signals to the risk engine. It also stores the KYC documents and supports their review by KYC officers.

\subsubsection{Core Banking}

\paragraph{Account.} The Account service maintains the customers and their accounts. It keeps one customer record for each person and any number of accounts of four products, which are savings, current, fixed deposit and recurring deposit. It records joint holders, account limits, freezes, nominees and saved payees. It applies the cooling period of new payees and delays any increase of a limit by 24 hours. It resolves pay identifiers to account numbers. It also keeps a hash-chained audit log of every change to an account, in which each entry contains a hash of the previous one. Any later change to the history therefore breaks the chain and becomes visible.

\paragraph{Transaction Intake.} The Transaction Intake service accepts every customer payment and runs it through the admission pipeline. It enforces idempotency, so a request that the client repeats is treated as the same payment. It calls the Rate Limit Coordinator, the Account service and the Risk Detection service in turn. Depending on the result, it posts the payment to the Ledger or asks the customer for a one-time password. It can also reserve the funds for an analyst or decline the payment. \Cref{sec:payment-lifecycle} describes this pipeline in detail.

\paragraph{Ledger.} The Ledger is the only service that moves money. It records every movement as a double-entry journal, in which the debits and credits of each journal sum to zero. It maintains the balances, places and releases holds on funds, and reverses journals when a payment is cancelled. It posts interest and runs a reconciliation that checks the sum of all balances against the journals. Database constraints reject an unbalanced journal, an overdrawn balance and any change to a posted entry.

\subsubsection{Admission Control and Risk}

\paragraph{Rate Limit Coordinator.} The Rate Limit Coordinator keeps the shared rate-limit state for all replicas of all services. It applies sliding-window limits to customers, accounts, devices and IP addresses for payments, and to identities and IP addresses for sign-in attempts. It stores the window counters in DynaStore as conflict-free counters, so every replica sees the same counts. It also hands out leases of permits, which let a busy caller admit several requests without a call per request. Operators change the policies through a configuration file in Git.

\paragraph{Risk Detection.} The Risk Detection service scores each payment against 21 weighted rules. The rules use features in DynaStore, such as the recent velocity and the usual payment amount of a customer. Other features record when a device or payee first appeared. They also use signals from other services, such as a sign-in from a new device or an early closure of a fixed deposit. The sum of the weights of the matching rules gives a score from 0 to 100. The score decides whether the payment proceeds, needs a one-time password, waits for an analyst or is blocked. New rules can run in shadow mode, in which they are evaluated and recorded but do not affect the decision.

\subsubsection{Fraud Response and Notification}

\paragraph{Fraud Alert.} The Fraud Alert service turns suspicious activity into work for the fraud team. It consumes transaction events, sign-in events and the signals of the stream processing job. It raises alerts, freezes accounts automatically when the evidence is strong, and groups related alerts into cases. Analysts work the cases through the staff console, and a second analyst must approve the unfreezing of an account. The service also writes short-lived signals and a deny list into DynaStore, which the risk engine reads on the next payment.

\paragraph{Notification.} The Notification Service delivers messages to customers. It runs as a Cloud Function that Pub/Sub triggers for each notification request. It always writes the message to the in-app inbox in Firestore. Depending on the type of message, it also sends an email through Gmail or Brevo and an RCS message through a simulated channel. Administrators can switch each channel and the email provider on or off at run time.

\subsubsection{Analytics}

\paragraph{Audit Analytics.} The Audit Analytics service is a Flink SQL job on Dataproc. It reads the stream of transaction events from Kafka and computes statistics over one-minute windows. These statistics include throughput, completed and declined payments, rate-limit decisions, risk decisions and the number of unique customers. It writes the statistics to a results topic. It also detects four fraud patterns over longer windows and writes them to a signals topic. \Cref{sec:streaming} describes the job.

\subsubsection{Deposits}

\paragraph{Deposits and End of Day.} The Deposits service manages fixed deposits and recurring deposits. It publishes the rate cards, quotes the maturity amount of a proposed deposit, and prices the early closure of an existing one. It also runs the end-of-day batch once a day. The batch accrues interest on savings accounts and deposits, collects recurring instalments, and pays out deposits that reach maturity. It records each business date once, so a repeated run has no further effect. The service also produces monthly statements and registers a fingerprint of each one, which lets anyone verify a statement file later.

\subsubsection{Leaderless Storage}

\paragraph{DynaStore.} DynaStore is the leaderless key-value store that implements the concept of the scenario. It runs as a StatefulSet of three replicas, each with its own persistent disk. Any replica can coordinate any request, and every key is stored on three replicas. The Rate Limit Coordinator, the Risk Detection service and the services that raise fraud signals use it through gRPC. \Cref{sec:concept} describes its design.

\begin{longtable}{@{}L{0.5cm}L{3.0cm}L{7.4cm}L{3.0cm}@{}}
\caption{Responsibilities, runtime and data ownership of each service.}\label{tab:services}\\
\toprule
No. & Service & Responsibility & Runtime and data \\
\midrule\endhead
1 & API Gateway & Single public entry point, token and session checks, role and ownership checks, PIN checks, web applications & GKE with autoscaling, no data \\
2 & Identity and Authentication & Registration, sign-in, tokens, sessions, devices, PINs, one-time passwords, KYC documents & GKE with autoscaling, Cloud SQL, Firestore, Cloud Storage \\
3 & Account & Customers, accounts, joint holders, limits, freezes, payees, pay identifiers, audit chain & GKE with autoscaling, Cloud SQL \\
4 & Transaction Intake & Admission pipeline, idempotency, challenges and holds & GKE with autoscaling, Cloud SQL \\
5 & Rate Limit Coordinator & Shared sliding-window limits and permit leases & GKE, DynaStore \\
6 & Ledger & Double-entry journals, balances, holds, reversals, interest, reconciliation & GKE with autoscaling, Cloud SQL \\
7 & Risk Detection & Scoring against 21 weighted rules, shadow rules & GKE with autoscaling, DynaStore \\
8 & Fraud Alert & Alerts, automatic freezes, cases, maker-checker review, deny list & GKE, Cloud SQL, Firestore \\
9 & Notification & In-app inbox, email and RCS delivery & Cloud Function, Firestore \\
10 & Audit Analytics & Windowed statistics and fraud pattern detection & Flink on Dataproc, Kafka, Cloud Storage \\
11 & Deposits and End of Day & Fixed and recurring deposits, end-of-day batch, statements & GKE, Cloud SQL, Cloud Storage \\
12 & DynaStore & Leaderless replicated key-value storage & GKE StatefulSet, persistent disks \\
\bottomrule
\end{longtable}

\subsection{Lifecycle of a Payment}
\label{sec:payment-lifecycle}

A payment shows how the services cooperate. \Cref{fig:payment-sequence} shows the sequence of calls, and this subsection explains each step.

The customer enters the payee, the amount and an optional note, and confirms the payment with the payment PIN. The web application generates an idempotency key for the attempt and sends the request to the API Gateway. The gateway checks the session and the PIN with the Identity service and forwards the request to the Transaction Intake service.

The Transaction Intake service first claims the idempotency key. If a payment with the same key already exists, the service returns the result of that payment and moves no further money. The service then asks the Rate Limit Coordinator over gRPC whether the customer, the account, the device and the IP address are within their limits. Next, it asks the Account service whether the customer holds an active account. It resolves the payee and applies the limit per payment and the limit for new payees.

The Transaction Intake service then sends the payment and its context to the Risk Detection service over gRPC. The risk engine returns a score, a decision and the reasons behind the decision. The decision leads to one of four paths.
\begin{itemize}
\item An allowed payment continues at once.
\item A challenged payment waits while the Identity service sends a one-time password to the customer. The payment continues when the customer enters the correct code, and it fails after a wrong code or a timeout.
\item A held payment reserves the funds through a hold in the Ledger. A fraud analyst then approves or rejects it, and the customer can also cancel it.
\item A blocked payment is declined, and no money moves.
\end{itemize}

For a payment that continues, the Transaction Intake service checks and reserves the daily limit under a lock. It then asks the Ledger over gRPC to post a double-entry journal. The Ledger uses its own idempotency key, so a retried request cannot post the same payment twice. Finally, the Transaction Intake service records the outcome and an outbox event in one database transaction. A relay publishes the event to Kafka, where the Fraud Alert service and the Flink job consume it. The Fraud Alert service then sends a receipt request through Pub/Sub to the Notification Service, which informs the customer.

The web application shows the outcome together with the six visible checks. These are the duplicate check, the rate limits, the account and payee check, the limits, the risk engine and the ledger posting. When a check stops the payment, the application shows which check stopped it and why. If a dependency fails during the pipeline, the payment stays in a processing state. The client can then retry with the same idempotency key, and the retry continues from where the first attempt stopped.

\mermaidplaceholder{payment-sequence}{Sequence of calls for a single payment.}{The client, gateway, Identity, Intake, Rate Limit, Account, Risk, Ledger, Kafka and the Notification function, with the PIN check, the idempotency check, the risk decision and its four outcomes, the double-entry posting, the outbox event to Kafka, and the receipt that Fraud Alert sends through Pub/Sub to the Notification function.}{9cm}

\subsection{Client Applications}
\label{sec:clients}

The API Gateway serves two web applications from the same origin. The customer application at the root path is a single-page application written in React and TypeScript. Vite compiles it into static files, which the repository stores in the static directory of the gateway. Its design targets phones first. It includes a swipe control for confirming payments, entry fields for PINs and one-time passwords, and an animated view of the checks of each payment. The original application, written in plain JavaScript, remains available under the path \code{/classic/}. It hosts the staff console and a few customer functions, such as statements and the camera QR scanner.

Serving both applications from the gateway has three advantages. The system needs only one public address. The browser makes every request to the same origin, so no cross-origin configuration is needed. The gateway can also apply a strict Content-Security-Policy, which allows scripts only from its own origin and forbids framing by other sites. Both applications keep the session in the storage of the browser tab, so closing the tab signs the user out.

% =====================================================================  5
\section{Interconnection Mechanisms}
\label{sec:interconnection}

This section describes how the services of \project{} communicate. It explains the four mechanisms that the system uses and the rule that assigns each interaction to a mechanism. It also describes the patterns that keep the communication reliable and secure.

\subsection{Overview}
\label{sec:comm-overview}

The services communicate through four mechanisms, which are REST, gRPC, Kafka and Pub/Sub. \Cref{fig:communication} shows which mechanism connects each pair of services. One rule decides the choice. Requests that a person makes and waits for use REST. Calls on the hot path of every payment use gRPC. Events that several consumers must receive in order, and possibly replay later, use Kafka. The trigger of the serverless function uses Pub/Sub. \Cref{tab:mechanisms} summarises the mechanisms and their uses.

\begin{table}[H]\centering\small
\caption{Communication mechanisms and their uses.}\label{tab:mechanisms}
\begin{tabularx}{\linewidth}{@{}L{2.0cm}L{2.6cm}YY@{}}
\toprule
Mechanism & Style & Used for & Reason for the choice \\
\midrule
REST & Synchronous request and response & Client to gateway, gateway to services, service calls off the payment path & Simple, widely supported and easy to inspect \\
gRPC & Synchronous remote procedure call & Intake to Rate Limit, Risk and Ledger, and every DynaStore call & Binary encoding, typed contracts, deadlines and status codes for frequent calls \\
Kafka & Asynchronous event stream & Transaction, sign-in, ledger and account events, fraud signals and analytics results & Durable, ordered per key and replayable, and the source and sink of the Flink job \\
Pub/Sub & Asynchronous message queue & Notification requests to the Cloud Function & Native trigger for Cloud Functions with automatic retries \\
\bottomrule
\end{tabularx}
\end{table}

\mermaidplaceholder{communication}{Communication mechanisms used between services.}{The same services as the decomposition figure, with REST, gRPC, Kafka and Pub/Sub links drawn in four different line styles and a legend.}{7cm}

\subsection{REST}
\label{sec:rest}

REST carries the requests of the web application and most calls between services. A REST interface exposes resources at URLs and uses the standard HTTP methods to read and change them. All REST messages in \project{} use JSON. Each service publishes an OpenAPI description of its interface, which FastAPI generates from the code.

The client sends every request to the API Gateway with a bearer token. The gateway verifies the token and the session, applies the checks of \cref{sec:responsibilities}, and forwards the request to the internal service. It sets an \code{X-Actor} header that names the authenticated caller and discards any such header that the client sent. Internal services read the identity of the caller only from this header. They never accept requests from outside the namespace, because the network policy of \cref{sec:network} blocks them.

Services also call one another over REST when the call is not on the payment path. For example, the Fraud Alert service freezes accounts through the Account service, and the Deposits service asks the Ledger to post interest. The Transaction Intake service also asks the Account service to authorise each payer. Every service-to-service call carries a timeout, so a slow dependency cannot hold a request open indefinitely.

\subsection{gRPC}
\label{sec:grpc}

gRPC carries the calls that occur on every payment. gRPC is a remote procedure call framework that sends Protocol Buffers messages over HTTP/2. Protocol Buffers is a compact binary format with a schema that defines every message and its fields. The repository keeps these schemas in the \code{proto} directory, and both the clients and the servers generate their code from them. A change that breaks a contract therefore fails at build time.

Four services offer gRPC interfaces. The Rate Limit Coordinator offers a \code{Check} call, and the Risk Detection service offers an \code{Evaluate} call. The Ledger offers calls to post journals, read balances, and create and release holds. DynaStore offers calls to read, write and increment values, together with the internal calls that replicas use among themselves.

gRPC suits these calls for three reasons. Binary messages are smaller and faster to encode than JSON. Every call carries a deadline, after which the client stops waiting. Every response carries a defined status code, which lets the caller tell a rejected request from an unavailable server.

gRPC keeps one long-lived HTTP/2 connection between a client and a server. A standard Kubernetes service would send all calls on that connection to a single pod, and pods added by the autoscaler would receive no traffic. The gRPC clients in \project{} therefore resolve a headless service, for which DNS returns the address of every pod. They then spread their calls across all pods in turn.

\subsection{Event Streaming with Kafka}
\label{sec:kafka}

Kafka carries the events that the system produces. Kafka is a distributed log, in which producers append events to topics and consumers read them in order at their own pace. Kafka keeps events after delivery, so a new consumer can read the history and a failed consumer can resume where it stopped.

The Kafka cluster runs inside GKE under the Strimzi operator, which manages Kafka through Kubernetes resources. The cluster runs in KRaft mode, in which Kafka manages its own metadata without a separate ZooKeeper ensemble. \Cref{tab:topics} lists the topics with their producers and consumers. Each topic has three partitions, and events with the same key, such as the same account, always go to the same partition. Each consumer therefore sees the events of one account in the order in which they happened.

\begin{table}[H]\centering\small
\caption{Kafka topics, their producers and their consumers.}\label{tab:topics}
\begin{tabularx}{\linewidth}{@{}L{3.6cm}L{3.2cm}Y@{}}
\toprule
Topic & Producer & Consumers \\
\midrule
\code{txn-events} & Transaction Intake & Audit Analytics (Flink), Fraud Alert \\
\code{auth-events} & Identity & Fraud Alert \\
\code{risk-signals} & Audit Analytics (Flink) & Fraud Alert \\
\code{analytics-results} & Audit Analytics (Flink) & Fraud Alert, which stores the statistics for the staff console \\
\code{ledger.journal-posted} & Ledger & Available to downstream consumers \\
\code{ledger.hold-changed} & Ledger & Available to downstream consumers \\
\code{account.events} & Account & Available to downstream consumers \\
\bottomrule
\end{tabularx}
\end{table}

The event stream decouples the payment path from analytics and fraud response. The Transaction Intake service publishes an event and returns its answer to the customer at once. It does not wait for the Flink job or the Fraud Alert service, and a slow or failed consumer cannot delay a payment. The same stream also lets the Flink job run on a separate managed cluster, because the job only needs network access to the brokers.

\subsection{Transactional Outbox}
\label{sec:outbox}

A service that changes its database and then publishes an event faces a dual-write problem. If the service crashes between the two steps, the database and the event stream disagree. A payment might complete without any event, or an event might announce a payment that never committed.

\project{} solves this problem with the transactional outbox pattern~\cite{outbox}, which \cref{fig:outbox} shows. A service never writes to Kafka directly. It inserts the event into an outbox table in the same database transaction as the business change. A relay runs as a sidecar container next to the service. It reads new rows from the outbox table and publishes them to Kafka, and it marks each row once Kafka acknowledges it. An event therefore exists exactly when the business change has committed.

The relay may publish an event twice if it crashes after Kafka acknowledges the event and before it marks the row. Delivery is therefore at least once. Every consumer tolerates duplicates by recording the identifiers of the events it has processed and skipping any event it has already seen. Several replicas of a service can run their relays at the same time. Each relay locks the rows it publishes and skips the rows that another relay holds.

\mermaidplaceholder{outbox}{Transactional outbox between a service database and Kafka.}{A service writes the business row and the outbox row in one transaction. The relay sidecar reads the outbox and publishes to Kafka, and consumers deduplicate by event identifier.}{4.5cm}

\subsection{Serverless Notification with Pub/Sub}
\label{sec:pubsub}

The Notification Service runs as a Cloud Function of the second generation. Cloud Functions runs code in response to events and charges only while the code runs. Pub/Sub, the managed message queue of Google Cloud, triggers the function through Eventarc for each message on the \code{notifications} topic.

Two services publish notification requests. The Identity service sends one-time passwords, security alerts and KYC decisions. The Fraud Alert service consumes the transaction events and sends payment receipts, declined-payment notices, fraud warnings and freeze notices. Each request names the customer, the type of message and the text. It may also name the channels to use. A one-time password, for example, goes to the inbox, by email and by RCS, while a receipt goes only to the inbox.

The function scales from zero to three instances and processes each message asynchronously. The publisher never waits for delivery. If the function fails, Pub/Sub retries the message until delivery succeeds. The function therefore treats every message as possibly repeated. It writes each inbox entry under the identifier of the message, so a retry overwrites the same entry. The function also reads its settings from Firestore. These settings let an administrator switch email or RCS off, or move email from Gmail to Brevo, without a new deployment.

% =====================================================================  6
\section{Transaction Admission and Fraud Detection}
\label{sec:fraud}

Transaction admission and fraud monitoring form the core of the scenario. This section first sets out the threats that \project{} defends against. It then describes each layer of defence, from the rate limiter that acts before any other work to the case management that acts minutes later. It closes with an evaluation in which a simulator attacks the running system.

\subsection{Threat Model}
\label{sec:threats}

The design considers the attacks that retail banks commonly face. \Cref{tab:threats} lists them with the property that makes each one hard to detect.

\begin{table}[htbp]\centering\small
\caption{Threats considered in the design.}\label{tab:threats}
\begin{tabularx}{\linewidth}{@{}L{3.4cm}Y@{}}
\toprule
Threat & Description \\
\midrule
Stolen password & An attacker obtains the password of a customer, for example through phishing, and signs in from a device of their own. \\
Account takeover & An attacker holds the password, the SIM card and the PIN of a customer. The attacker can then receive one-time passwords and pass every credential check. \\
Payment bots & A script submits payments at machine speed, for example to empty an account before the customer notices. \\
Card testing & An attacker sends tiny payments to learn whether an account works, and then sends a large payment. \\
Structuring & A customer splits a large sum into payments that stay just below the threshold of Rs 50,000, which would otherwise attract scrutiny. \\
Mule rings & Many victims are tricked into paying one collection account. Each payment looks ordinary on its own. \\
Account changes before theft & An attacker changes the email address or phone number of a customer, or resets a credential, before moving money. \\
Staff misuse & A single member of staff tries to release a frozen account or approve a payment without oversight. \\
\bottomrule
\end{tabularx}
\end{table}

These threats differ in what they reveal. A bot reveals itself through volume, and a takeover through an unfamiliar device and place. Card testing and structuring reveal themselves only through a short sequence of payments from one account. A mule ring reveals itself only when the payments of many customers are viewed together. No single check can see all of these signs, which motivates the layered design that follows.

\subsection{Defence in Depth}
\label{sec:depth}

Defence in depth places several independent controls in the path of an attack, so that an attack that passes one control meets the next. \project{} applies five layers to every payment. \Cref{fig:fraud-layers} shows the path of a payment through them.

The layers differ in speed and in breadth of view. Admission control is the cheapest layer and sees only counts. The authorisation layer sees the account and its limits. The risk engine sees the recent history of the customer, the device and the payee, and it decides within tens of milliseconds. Stream analytics sees all payments over windows of up to an hour, but it reports after the window closes. Case management brings in human judgement. The slow layers feed their findings back into the fast layers through signals in DynaStore. The system therefore improves its real-time decisions with knowledge that only the slow layers can gather.

\mermaidplaceholder{fraud-layers}{Path of a payment through the five layers of defence and the decisions each layer can take.}{A payment enters admission control, then authorisation and policy, then risk scoring with its four score bands. Accepted payments reach the ledger and the event stream, where stream analytics and the fraud service act after the fact and feed signals back.}{7cm}

\subsection{Admission Control}
\label{sec:admission}

The Rate Limit Coordinator decides whether a request may proceed at all. It applies limits to several scopes at once, and a request must pass every scope. \Cref{tab:ratelimits} lists the policies. Each policy allows at most a given number of permits within a sliding window.

\begin{table}[H]\centering\small
\caption{Rate-limit policies for payments and sign-in attempts.}\label{tab:ratelimits}
\begin{tabularx}{\linewidth}{@{}L{2.4cm}L{2.6cm}Y@{}}
\toprule
Action & Scope & Limit \\
\midrule
Payment & Customer & 10 per minute and 100 per hour \\
Payment & Account & 10 per minute \\
Payment & Device & 20 per minute \\
Payment & IP address & 120 per minute, with leases of 10 permits \\
Sign-in & Email or phone & 5 per 5 minutes \\
Sign-in & IP address & 30 per minute \\
\bottomrule
\end{tabularx}
\end{table}

The scopes complement one another. The customer and account limits stop a bot that uses one stolen login. The device limit stops a single machine that cycles through many accounts. The IP limit stops a burst from one network. The sign-in limit on each email address or phone number slows the guessing of one customer's password. The sign-in limit on each IP address slows guessing across many customers.

The coordinator keeps the window counters in DynaStore. Every replica of every service therefore sees the same counts, and an attacker gains nothing by spreading requests across replicas. The counters are conflict-free replicated counters, which \cref{sec:concept} describes. Two replicas can increment the same counter at the same moment without losing either increment. For the busiest scope, the IP address, a coordinator replica takes a lease of ten permits at a time and hands them out locally. Leases trade a small loss of precision for far fewer calls to the store.

The coordinator must also decide what happens when DynaStore is unreachable. Refusing every payment would turn a storage fault into an outage of the bank. Admitting every payment would remove the protection exactly when an attacker might exploit a fault. The coordinator therefore admits payments of up to Rs 2,000 and refuses larger ones. Small payments fail open, and large payments fail closed.

\subsection{Identity-Based Controls}
\label{sec:identity-controls}

The second group of controls binds every payment to a verified person on a known device. These controls act before the risk engine and also feed it.

\paragraph{Device binding.} The application identifies each browser by a random device identifier. When a customer signs in on a device that the bank has not seen before, the correct password is not enough. The Identity service sends a six-digit code to the email address and to the inbox of the customer on their other devices. The session starts only after the customer enters this code. A stolen password alone therefore opens nothing.

\paragraph{Two PINs.} Each customer holds a six-digit payment PIN and a four-digit app PIN. The payment PIN approves payments, new payees, higher limits and joint holders. The app PIN unlocks the application and approves deposits and nominees. The API Gateway verifies the required PIN with the Identity service before the request reaches the service that does the work. Three wrong attempts lock a PIN for 24 hours, and the customer receives an alert. The Identity service stores each PIN as an HMAC-SHA256 digest with a salt and a secret pepper. A copy of the database therefore does not reveal the PINs.

\paragraph{Sessions.} Access tokens live for 15 minutes, but the gateway also checks the session behind every token with the Identity service. Signing out, removing a device or blocking all access therefore takes effect at once. A customer session locks after five minutes without activity and unlocks with the app PIN. A staff session ends after 15 minutes without activity. Each refresh token works only once. If an old refresh token appears again, someone has copied it, so the Identity service ends the session and raises a signal.

\paragraph{Delays on risky changes.} Lowering a payment limit takes effect at once, but raising one takes effect only after 24 hours. A new payee can receive at most Rs 10,000 per payment during the first 24 hours. These delays give the customer time to notice a change that someone else made.

\paragraph{Signals from account changes.} The Identity service writes short-lived signals into DynaStore after certain account events. These events are a change of email address or phone number, a reset of a password or PIN, a locked PIN and a replayed refresh token. It also writes signals at sign-in for a new device and for impossible travel. Impossible travel means that the distance between two consecutive sign-ins, divided by the time between them, exceeds 900 kilometres per hour. The risk engine weighs all of these signals in its next decision.

\subsection{Real-Time Risk Scoring}
\label{sec:risk}

The Risk Detection service scores every payment that passes admission and authorisation. It evaluates the 21 rules of \cref{tab:rules}. Every rule that matches adds its weight to the score, and the total is capped at 100.

\begin{longtable}{@{}L{4.4cm}L{1.1cm}L{8.6cm}@{}}
\caption{Rules of the real-time risk engine and their weights.}\label{tab:rules}\\
\toprule
Rule & Weight & Condition \\
\midrule\endhead
\code{DENY\_LIST} & 100 & The payee or the device received confirmed fraud in the last 30 days. \\
\code{VELOCITY\_1M} & 35 & More than 5 payment attempts from the customer in one minute. \\
\code{VELOCITY\_1H} & 20 & More than 20 payment attempts from the customer in one hour. \\
\code{AMOUNT\_ANOMALY} & 30 & At least Rs 5,000 and more than five times the usual amount, given at least five earlier payments. \\
\code{HIGH\_AMOUNT} & 10 & At least Rs 40,000. \\
\code{NEW\_DEVICE} & 20 & The device first appeared less than 24 hours ago. \\
\code{NEW\_PAYEE} & 10 & At least Rs 5,000 to a payee first paid less than 24 hours ago. \\
\code{ATO\_PATTERN} & 40 & A new device, a new payee and at least Rs 10,000 together. \\
\code{MICRO\_PROBING} & 60 & At least Rs 5,000 after three or more recent payments of at most Rs 10. \\
\code{DORMANT\_ACCOUNT} & 55 & At least Rs 5,000 from an account idle for 90 days. \\
\code{ODD\_HOURS} & 10 & Between 1 and 5 in the morning, India time. \\
\code{STRUCTURING} & 60 & The third recent payment within 5 percent below Rs 50,000. \\
\code{SESSION\_NEW\_DEVICE\_LOGIN} & 10 & The session began with a new-device sign-in. \\
\code{IMPOSSIBLE\_TRAVEL} & 40 & The session began with an impossible-travel sign-in. \\
\code{FD\_BREAK\_AND\_FLEE} & 35 & A fixed deposit was closed early, and the money goes to an untrusted payee. \\
\code{PATTERN\_SIGNAL\_SENDER} & 30 & Stream analytics flagged the paying account. \\
\code{MULE\_PAYEE} & 60 & Stream analytics flagged the receiving account as a mule. \\
\code{CONTACT\_CHANGED} & 20 & The email address or phone number changed in the last three days. \\
\code{CREDENTIAL\_RESET} & 20 & A password or PIN was reset or changed in the last day. \\
\code{PIN\_LOCKED} & 15 & A PIN locked after three wrong attempts in the last day. \\
\code{TOKEN\_REUSE} & 40 & A replayed refresh token was detected in the last day. \\
\bottomrule
\end{longtable}

The rules draw on three kinds of input. Velocity counters and behavioural features, such as the usual amount and the first appearance of a device or payee, live in DynaStore. The Risk Detection service updates them with every payment. Signals from other services also live in DynaStore and expire on their own. The remaining inputs, such as the amount and the time of day, come with the request itself.

The score falls into one of four bands, which decide the outcome. A score below 30 allows the payment. A score from 30 to 59 challenges it, so the customer must confirm the payment with a one-time password. A score from 60 to 94 holds it. The Ledger then reserves the funds, and a fraud analyst approves or rejects the payment. A score of 95 or more blocks it.

The band boundaries follow from the weights. Only the deny list, or several strong signals together, can reach the blocking band. A weak signal, such as the time of day, can raise a payment to a challenge or a hold but never to a block. Holding a doubtful payment protects the customer without the cost of refusing a genuine payment outright. Signals about account changes count only when the payee is untrusted. A customer who changes a phone number can therefore still pay their landlord without friction.

The risk engine also handles its own failures. If DynaStore is unreachable, the velocity and history rules cannot run. The engine then holds every payment of Rs 20,000 or more for review. The rules live in a configuration file in Git, and the service reloads them without a restart. A new rule can run in shadow mode first, in which the engine evaluates and records it without adding its weight. The team can then measure how often a rule would fire before it affects customers. Every decision returns the names of the matching rules, and the application shows them to the customer in plain language.

\subsection{Stream Pattern Detection}
\label{sec:patterns}

Some fraud is invisible in any single payment. The Audit Analytics job on Flink therefore examines the stream of transaction events over time windows. A tumbling window divides time into fixed, non-overlapping intervals. A hopping window has a fixed length but starts at regular steps that are shorter than its length, so consecutive windows overlap. \Cref{tab:patterns} lists the signals that the job emits.

\begin{table}[H]\centering\small
\caption{Fraud signals detected by stream analytics.}\label{tab:patterns}
\begin{tabularx}{\linewidth}{@{}L{3.2cm}L{3.6cm}Y@{}}
\toprule
Signal & Window & Condition \\
\midrule
\code{VELOCITY\_BURST} & 1 minute, tumbling & At least 8 payment attempts from one account \\
\code{DECLINE\_BURST} & 1 minute, tumbling & At least 5 declined attempts from one account \\
\code{MICRO\_PROBING} & 5 minutes, hopping every minute & At least 3 completed payments of at most Rs 10 from one account \\
\code{STRUCTURING} & 60 minutes, hopping every 5 minutes & At least 3 payments from one account between Rs 47,500 and Rs 49,999.99 \\
\code{MULE\_FAN\_IN} & 10 minutes, hopping every minute & At least 5 distinct senders paying at least Rs 5,000 each into one account \\
\bottomrule
\end{tabularx}
\end{table}

The bursts and the probing signal confirm, over a wider window, what the real-time rules suspect within seconds. The structuring signal catches a customer who spaces payments out to avoid the real-time rule. The mule fan-in signal is the clearest case for stream analytics. Each payment into the mule account comes from a different customer, who has an ordinary history and pays an ordinary amount. The risk engine of each payer therefore sees nothing unusual. Only a view across customers, grouped by the receiving account, reveals the pattern. The minimum amount of Rs 5,000 keeps ordinary shops, which receive many small payments, out of this signal.

The job writes each signal to the \code{risk-signals} topic. The Fraud Alert service consumes the signal and writes a short-lived marker into DynaStore. The next payment to a flagged mule account matches the \code{MULE\_PAYEE} rule, and the next payment from a flagged sender matches the \code{PATTERN\_SIGNAL\_SENDER} rule. In this way, stream analytics changes the real-time decision on the very next payment.

\subsection{Response and Case Management}
\label{sec:response}

The Fraud Alert service turns events and signals into protective action and into work for the fraud team. Its behaviour depends on the event.

A held payment opens a case for an analyst, with high severity when the risk score is 75 or more. The customer receives a message that asks whether they made the payment. A blocked payment raises a high-severity alert and informs the customer. A customer who answers that they did not make a held payment triggers the strongest response. The payment is cancelled, the account is frozen, and an account takeover case opens.

A stream signal raises an alert. Because hopping windows report the same pattern repeatedly, the service keeps one signal of each type for each account within ten minutes. It treats a signal as serious when the signal is a mule fan-in or when two different signals concern the same account within ten minutes. A serious signal freezes the account automatically, opens a high-severity case and informs the account holder.

When an analyst confirms fraud, the service cancels the held payment and freezes the account. It also adds the payee and the device to a deny list in DynaStore for 30 days. Any later payment to that payee or from that device then matches the \code{DENY\_LIST} rule and is blocked. Unfreezing an account follows the maker-checker rule. One analyst requests the release, and a different analyst must approve it.

Customers hold three controls of their own. They can cancel a payment that the bank has held. They can report a completed payment, which opens a case. They can also block everything in one action. This action freezes every account of the customer, signs out every device and opens a case for a fraud officer. Incoming payments still arrive while the accounts are frozen.

\subsection{Evaluation with a Fraud Simulator}
\label{sec:fraud-eval}

A fraud simulator in the repository tests the layers against the deployed system. It creates fresh customers and runs eight scenarios through the public API. The first scenario is an ordinary customer, and the other seven are attacks. \Cref{tab:simulator} lists the outcome of each scenario in the final run against the cloud deployment.

\begin{table}[H]\centering\small
\caption{Outcome of the eight fraud simulator scenarios against the cloud deployment.}\label{tab:simulator}
\begin{tabularx}{\linewidth}{@{}L{3.1cm}L{3.6cm}Y@{}}
\toprule
Scenario & Layer that responded & Outcome \\
\midrule
Ordinary customer & None & Three small payments completed without friction \\
Payment bot & Rate Limit Coordinator and velocity rule & Of 30 rapid payments, 5 completed, 5 needed a one-time password and 20 were rate-limited \\
Stolen password & Identity, device binding & The sign-in from a new device stopped at the code sent to the customer \\
Account takeover & Risk Detection & Score 100 from a new device, impossible travel, a new payee and the takeover pattern, so the payment was blocked \\
Micro-payment probing & Risk Detection & Three payments of Rs 2 completed, and the following Rs 9,500 payment was held with score 70 \\
Structuring & Risk Detection & The first two payments just below Rs 50,000 completed, and the third and fourth were held with score 70 \\
Fixed deposit break and transfer & Risk Detection with a signal from Deposits & The transfer to a new payee needed a one-time password, with score 45 \\
Mule ring & Audit Analytics and Fraud Alert & Six senders within ten minutes raised a mule fan-in signal, which froze the account and opened a case \\
\bottomrule
\end{tabularx}
\end{table}

The results show that each attack meets the layer designed for it. The ordinary customer passed every check without a challenge, which shows that the layers do not burden normal use. The mule ring passed every real-time check, as expected, and only stream analytics caught it. This outcome confirms the value of the slow layer and of its feedback into the fast layers.

The evaluation has two limits. First, the team chose the rule weights and thresholds by reasoning about each threat, and it tested them only on synthetic traffic. A bank would tune them against labelled historical data and track the rates of false alarms and missed fraud. Second, the simulator tests one instance of each attack. It shows that the layers work as designed, but it does not measure how well they generalise to attacks that differ from these scenarios.

% =====================================================================  7
\section{Data Architecture}
\label{sec:data}
\plan{This section describes the three kinds of data in the system and the store chosen for each. It covers the storage products, the leaderless store that implements the assigned concept, and the stream processing job that turns the event stream into statistics and signals.}
\subsection{Storage Products and Data Ownership}
\label{sec:storage}
\plan{This section describes each storage product, the data it holds and the reason for the choice. It covers Cloud SQL for PostgreSQL, Firestore, Cloud Storage, DynaStore and Kafka. It also explains how schema migrations follow Git through the pre-sync job, and how KYC documents and statements are stored and protected.}

\begin{table}[H]\centering\small
\caption{Storage products, the data each one holds, and the reason for the choice.}\label{tab:storage}
\begin{tabularx}{\linewidth}{@{}L{3cm}YY@{}}
\toprule
Product & Data held & Reason for the choice \\
\midrule
Cloud SQL (PostgreSQL) & & \\
Firestore & & \\
Cloud Storage & & \\
DynaStore & & \\
Kafka & & \\
\bottomrule
\end{tabularx}
\end{table}

\subsection{Leaderless Replication in DynaStore}
\label{sec:concept}

\subsubsection{Motivation}
\plan{This subsection explains why rate-limit counters and fraud features suit leaderless replication. They are written on every request, they tolerate slight staleness, and they must survive node failures.}

\subsubsection{Partitioning with Consistent Hashing}
\plan{This subsection describes the hash ring with virtual nodes and the preference list of each key.}

\mermaidplaceholder{dynastore-ring}{Consistent hashing ring of the three DynaStore nodes and the preference list of one key.}{A ring with virtual node positions for three nodes and one key placed on it, with arrows to the first three distinct nodes clockwise.}{6cm}

\subsubsection{Replication and Tunable Quorums}
\plan{This subsection describes replication with N equal to 3, and the read and write quorums R and W that each request chooses. It explains why R plus W greater than N gives a read that overlaps the latest write.}

\subsubsection{Versioning and Conflict Resolution}
\plan{This subsection describes vector clocks for registers and PN-counter CRDTs for counters, and how each one merges concurrent updates without losing data.}

\subsubsection{Failure Handling}
\plan{This subsection describes the sloppy quorum, hinted handoff, read repair, anti-entropy and failure detection. It also describes the fault injection endpoint that the demonstration uses to cut a node off.}

\subsubsection{Application in the Bank and Its Boundaries}
\plan{This subsection lists the keyspaces that the bank stores in DynaStore with their writers and consistency levels. It then explains why balances are kept in PostgreSQL instead, because concurrent debits accepted by two coordinators would amount to a double spend. It closes with the degraded behaviour when DynaStore is unreachable.}

\subsubsection{Observed Behaviour}
\plan{This subsection reports what the team observed during the load test and the partition demonstration. Hints accumulated while nodes were overloaded and drained once the load stopped, which shows the store recovering without intervention.}

\screenshotplaceholder{dynastore-ui}{DynaStore status page showing ring ownership, key counts and pending hints.}{The status page during or after the load test.}{4.5cm}

\subsection{Stream Processing on Dataproc}
\label{sec:streaming}

\subsubsection{Job Design}
\plan{This subsection describes the Flink SQL job that reads the transaction topic with event-time watermarks. It lists the tumbling and hopping windows in a table together with the outputs and the topics they feed. It states that the per-minute window counts unique customers.}

\mermaidplaceholder{flink-pipeline}{Stream processing pipeline from the transaction topic to the results and signal topics.}{The txn-events topic feeds the Flink job on Dataproc. Windowed aggregations write to analytics-results and risk-signals, Fraud Alert consumes the signals, and checkpoints go to Cloud Storage.}{5.5cm}


\subsubsection{Deployment on Dataproc}
\plan{This subsection describes the single-node Dataproc cluster, the YARN session and the connection to Kafka inside the VPC. It records the engineering problems that the team solved, which were the discovery of the Kafka nodes, the default parallelism, a missing library of the Kafka connector, and the stalling of event-time windows when input stops.}

\subsubsection{Results}
\plan{This subsection reports the output of the job for a simulated mule ring and shows the resulting fraud case.}

\screenshotplaceholder{flink-results}{Window statistics and fraud signals produced by the stream processing job.}{The analytics-results and risk-signals output, the Stream analytics table and the open fraud cases.}{5cm}


% =====================================================================  8
\section{Quality Attributes and Evaluation}
\label{sec:quality}
\plan{This section evaluates the system against the qualities that a bank needs, which are scalability, resilience, observability and security.}
\subsection{Scalability and Load Testing}
\label{sec:scalability}

\subsubsection{Autoscaling Configuration}
\plan{This subsection lists the Horizontal Pod Autoscalers with their targets and replica limits in a table. It explains why the replica limits of the database services follow the connection budget of Cloud SQL.}

\subsubsection{Load Test Methodology}
\plan{This subsection describes the k6 test, its two scenarios for reads and payments, its ramp profile and its pass criteria for error rate and latency.}

\subsubsection{Results and Analysis}
\plan{This subsection reports the three load test runs in a table and analyses each one. The first run exposed a deadlock in the database connection handling, which the team diagnosed from logs and fixed. The second run showed that node CPU had become the limit. The third run, on three nodes, met both criteria. Grafana figures show the request rate, the error rate and the autoscaling of the services.}

\begin{table}[H]\centering\small
\caption{Results of the three load test runs.}\label{tab:loadtest}
\begin{tabular}{@{}lllll@{}}
\toprule
Run & Configuration & Failed requests & p95 latency & Payments completed \\
\midrule
1 & Before the connection fix, 2 nodes & 63.7\% & 20.0 s & 487 \\
2 & After the fix, 2 nodes & 4.2\% & 11.7 s & 3,973 \\
3 & After the fix, 3 nodes & 0.01\% & 1.4 s & 4,529 \\
\bottomrule
\end{tabular}
\end{table}

\screenshotplaceholder{grafana-loadtest}{Grafana dashboard during the final load test.}{Requests per second by service, errors, latency and HPA replicas during run 3.}{6cm}



\subsection{Resilience}
\label{sec:resilience}
\plan{This subsection describes the liveness and readiness probes, the idempotency keys on every write path, the retries, the disruption budget of the edge, the tolerance of Spot node loss, and the retrying end-of-day batch.}
\subsection{Observability}
\label{sec:observability}

\subsubsection{Metrics and Dashboards}
\plan{This subsection describes the metrics that every service exposes, the Prometheus deployment, and the panels of the Grafana dashboard for service health, cluster health and domain behaviour.}

\subsubsection{Centralised Logging}
\plan{This subsection describes the structured JSON logs, their collection by Cloud Logging from GKE and from the Dataproc job, and an example of how the team used them to locate the connection deadlock.}

\screenshotplaceholder{cloud-logging}{Cloud Logging view of the service logs.}{A Logs Explorer query across the bank namespace.}{4.5cm}

\subsection{Security}
\label{sec:security}
\plan{This section describes authentication with sessions, two PINs and new-device codes, the authorisation model for customers and staff, the handling of secrets, the network controls, the protection of personal data, and the tamper-evident audit log.}


% =====================================================================  9
\section{Design Rationale}
\label{sec:rationale}
\plan{This section collects the main design decisions in one place, which a design document must justify. Each decision states the choice, the alternatives considered and the reason for the choice. The decisions cover the cloud provider, GKE with Spot nodes, Kafka inside the cluster, the split between PostgreSQL and DynaStore, gRPC on the hot path, the outbox pattern, the Cloud Function for notifications, Flink SQL on Dataproc, Terraform layering, pull-request releases with ArgoCD, and the edge with Traefik.}

\begin{longtable}{@{}L{3.2cm}L{4.3cm}L{6.6cm}@{}}
\caption{Principal design decisions, the alternatives considered and the reasons for each choice.}\label{tab:rationale}\\
\toprule
Decision & Alternatives considered & Reason \\
\midrule\endhead
Google Cloud Platform & & \\
GKE zonal cluster on Spot nodes & & \\
Kafka on Strimzi inside GKE & & \\
Balances in PostgreSQL, counters in DynaStore & & \\
gRPC for hot-path calls & & \\
Transactional outbox & & \\
Cloud Function for notifications & & \\
Flink SQL on Dataproc & & \\
Three Terraform layers & & \\
Release by deploy pull request & & \\
Traefik behind a Google load balancer & & \\
\bottomrule
\end{longtable}


% =====================================================================  10
\section{Conclusion and Future Work}
\label{sec:conclusion}

\subsection{Limitations and Future Work}
\label{sec:limitations}
\plan{This subsection states the known limitations, such as the single zonal database, the simulated RCS channel, the dependence of the Flink job on node addresses, the uneven balancing of long-lived gRPC connections, and the hand-tuned risk weights. It suggests a concrete improvement for each.}

\subsection{Conclusion}
\plan{This subsection restates the problem, the architecture and the evidence that the system meets each requirement, in two short paragraphs.}

% =====================================================================  references
\begin{thebibliography}{9}
\bibitem{dynamo} G. DeCandia, D. Hastorun, M. Jampani, et al. Dynamo: Amazon's highly available key-value store. In \emph{Proceedings of the 21st ACM Symposium on Operating Systems Principles}, 2007.
\bibitem{ddia} M. Kleppmann. \emph{Designing Data-Intensive Applications}. O'Reilly Media, 2017.
\bibitem{crdt} M. Shapiro, N. Pregui\c{c}a, C. Baquero, and M. Zawirski. Conflict-free replicated data types. In \emph{Stabilization, Safety, and Security of Distributed Systems}, 2011.
\bibitem{flink} P. Carbone, A. Katsifodimos, S. Ewen, et al. Apache Flink: Stream and batch processing in a single engine. \emph{IEEE Data Engineering Bulletin}, 38(4), 2015.
\bibitem{outbox} C. Richardson. \emph{Microservices Patterns}. Manning Publications, 2018.
\bibitem{k8s} The Kubernetes Authors. Horizontal Pod Autoscaling. Kubernetes documentation, 2026.
\bibitem{argocd} Argo Project. Argo CD declarative GitOps continuous delivery. Project documentation, 2026.
\end{thebibliography}

% =====================================================================  appendices
\appendix

\section{Repository Structure}
\plan{This appendix gives the top-level layout of the repository with one sentence for each directory.}

\section{Deployment Procedure}
\plan{This appendix lists the commands that bring the environment up and down and the order in which to run them.}

\section{Task Division}
\plan{This appendix reproduces the content of task\_division.txt with the ID number, name, requirements and services of each member.}

\end{document}
